\documentclass[11pt,a4paper]{article}
\usepackage[margin=25mm]{geometry}
\usepackage[T1]{fontenc}
\usepackage{amsmath,amssymb,amsthm,mathtools}
\usepackage{newtxtext,newtxmath}
\usepackage{microtype,bm,booktabs,array}
\usepackage{xcolor}
\definecolor{LinkBlue}{RGB}{35,58,115}
\usepackage[colorlinks,linkcolor=LinkBlue,citecolor=LinkBlue,urlcolor=LinkBlue]{hyperref}
\hypersetup{pdftitle={Contraction, time reversal, and positive bosonic traces},pdfauthor={}}
\newcommand{\Tr}{\operatorname{Tr}}
\newcommand{\tr}{\operatorname{tr}}
\newcommand{\diag}{\operatorname{diag}}
\newcommand{\ee}{\mathrm e}
\newcommand{\ii}{\mathrm i}
\newcommand{\trans}{\mathsf T}
\newcommand{\mc}[1]{\mathcal{#1}}
\newcommand{\bb}[1]{\bm{#1}}
\newcommand{\hc}{\dagger}
\newtheorem{theorem}{Theorem}
\numberwithin{equation}{section}
\allowdisplaybreaks[1]
\setlength{\parskip}{3pt}
\setlength{\emergencystretch}{2em}

\title{Contraction, time reversal, and positive bosonic traces}
\author{}
\date{September 16, 2026}

\begin{document}
\maketitle
\vspace{-1.5em}
\begin{abstract}
Bosonic auxiliary-field weights require control of both convergence and
sign. Damping suppresses high occupations and gives a contraction inequality
for the finite propagation matrix. A common physical time reversal makes
the trace real. We show that damped quadratic evolution can be connected
continuously to independent harmonic oscillators while its trace stays finite
and nonzero. The trace is therefore positive for either square, $+I$ or
$-I$, of time reversal. This gives a direct interpretation of the two
standard bosonic contraction semigroups. We also place reflection-positive
Gaussian products in this framework by writing their physical metric and
an explicit change of boson modes. A paired auxiliary-field example
illustrates strict damping combined with unitary field factors.
A state-evolution argument also permits amplifying slices.
\end{abstract}

\section{The physical question}

After an auxiliary field has been fixed, a bosonic quantum Monte Carlo
calculation encounters an ordered product
\begin{equation}
 \hat G=\ee^{\hat h_1}\ee^{\hat h_2}\cdots\ee^{\hat h_L},
 \qquad W=\Tr\hat G.
 \label{eq:weight}
\end{equation}
There are finitely many modes, but each mode has arbitrarily large occupation.
Each $\hat h_\ell$ is quadratic in creation and annihilation operators. Its
coefficient includes the time step and the sign of the imaginary-time
evolution. Pair creation and annihilation are allowed, and the generators may
be non-Hermitian. The trace is over the full Fock space, without a projection
onto a chosen particle number or parity.
Such quadratic evolutions are called Gaussian because they preserve the
Gaussian form of oscillator wavefunctions whenever the evolved state exists.

For one oscillator, the distinction between convergence and sign is already
visible. With $\hat n=\hat a^+\hat a$,
\begin{equation}
 \Tr\ee^{-\beta\hat n}
 =\sum_{n=0}^{\infty}\ee^{-\beta n}
 =\frac{1}{1-\ee^{-\beta}}>0,\qquad \beta>0.
 \label{eq:thermal_number}
\end{equation}
At $\beta<0$ the occupation sum diverges, although the expression on the far
right is finite and negative. The present note concerns the ordinary
convergent trace. It does not assign that negative continuation to a thermal
sum.

For non-Hermitian products, convergence will mean that $\hat G$ has a finite
sum of singular values, the square roots of the eigenvalues of
$\hat G^\dagger\hat G$. This is the condition usually called \emph{trace
class}. It makes $W$ independent of the orthonormal basis used to compute it
and permits absolutely convergent occupation sums. In physical terms, large
occupations must be sufficiently suppressed in the final evolution.

Time reversal is an antiunitary operator $\Theta$. It conjugates numerical
coefficients, preserves the order of products, and takes a creation operator
to the adjoint of the transformed annihilation operator. We require the
\emph{same} symmetry for all slices,
\begin{equation}
 \Theta\hat h_\ell\Theta^{-1}=\hat h_\ell
 \quad\text{for every }\ell.
 \label{eq:slice_trs}
\end{equation}
This implies $\Theta\hat G\Theta^{-1}=\hat G$ and $W=W^*$. Reality is the
first step. The central task is to show why the real number $W$ is positive.

We first solve this task without pairing. With pairing, the useful physical
picture is damping combined with unitary motion. Section~\ref{sec:contraction}
translates damping into the contraction inequality, and Sec.~\ref{sec:trace}
uses it to prove positivity by a continuous connection to an elementary
harmonic-oscillator partition function. Sections~\ref{sec:standard} and
\ref{sec:rp} identify the two
standard forms and the place of reflection positivity within them.

\section{Number conservation: convergence makes one time reversal sufficient}
\label{sec:number}

Consider
\begin{equation}
 \hat h_\ell=\xi_\ell+
 \sum_{ij}\hat a_i^+(K_\ell)_{ij}\hat a_j,
 \qquad B=\ee^{K_1}\cdots\ee^{K_L}.
 \label{eq:number_generator}
\end{equation}
Time reversal makes each scalar $\xi_\ell$ real. Their exponentials are
positive, so we temporarily omit them. On the one-particle sector $\hat G$
acts as $B$. On the $n$-particle sector it acts as $n$ copies of $B$ restricted
to symmetric bosonic states. Consequently, if the eigenvalues of $B$ are
$\lambda_1,\ldots,\lambda_N$, the occupation sum gives
\begin{equation}
 W=\prod_{j=1}^{N}\sum_{n_j=0}^{\infty}\lambda_j^{n_j}
   =\frac{1}{\det(I-B)}.
 \label{eq:number_trace}
\end{equation}
The formula also holds when $B$ is not diagonalizable, by continuity from
diagonalizable matrices in the convergent region.

Convergence constrains all $|\lambda_j|$ to be smaller than one. A direct
way to see the stronger operator condition is to apply the singular-value
decomposition $B=V\diag(s_1,\ldots,s_N)W^\dagger$ on every bosonic sector.
The two unitary mode rotations do not change singular values. Summing the
remaining singular values of the many-body operator gives
\begin{equation}
 \Tr\sqrt{\hat G^\dagger\hat G}
 =\prod_j\sum_{n_j=0}^{\infty}s_j^{n_j}.
 \label{eq:number_singular}
\end{equation}
Thus a genuine convergent trace requires $s_j<1$ for every $j$, which in
particular gives $|\lambda_j|<1$. Merely locating the eigenvalues inside the
unit circle is insufficient to establish convergence for a nonnormal $B$.

In a basis of annihilation operators, write time reversal as
\begin{equation}
 \Theta\hat a\Theta^{-1}=U\hat a,
 \qquad U^\dagger U=I,\qquad \overline U U=\pm I.
 \label{eq:trs_a}
\end{equation}
Slice invariance gives $U^\dagger\overline B U=B$. Hence the eigenvalues of
$B$ are real or occur in complex-conjugate pairs. Grouping the factors in
Eq.~\eqref{eq:number_trace} accordingly yields
\begin{equation}
 \det(I-B)=
 \prod_{\lambda\in\mathbb R}(1-\lambda)
 \prod_{\operatorname{Im}\lambda>0}|1-\lambda|^2>0.
 \label{eq:number_positive}
\end{equation}
Every real factor is positive because $-1<\lambda<1$. No degeneracy of the
real eigenvalues is needed. Therefore either square in
Eq.~\eqref{eq:trs_a} is sufficient for $W>0$ when the trace converges.

For $\Theta^2=-I$ on the one-particle operators, Kramers pairing remains a
useful alternative argument. Equation~\eqref{eq:number_positive} explains
why convergence also permits $\Theta^2=+I$. It supplies the positive sign of
each unpaired real factor.

\section{Damping and the contraction semigroup}
\label{sec:contraction}

With pairing, the growing and decaying occupation sectors are mixed.
Positions and momenta provide a simple way to identify what damps the
evolution. Introduce Hermitian coordinates
\begin{equation}
 \hat a_j=\frac{\hat x_j+\ii\hat p_j}{\sqrt2},\qquad
 \hat a_j^+=\frac{\hat x_j-\ii\hat p_j}{\sqrt2},\qquad
 \hat r=\binom{\hat x}{\hat p},\qquad
 [\hat r_\mu,\hat r_\nu]=\ii\Omega_{\mu\nu},\quad
 \Omega=\begin{pmatrix}0&I\\-I&0\end{pmatrix}.
 \label{eq:quadratures}
\end{equation}
Apart from a real scalar, write
\begin{equation}
 \hat h=\frac12\hat r^\trans H\hat r,\qquad
 H^\trans=H,\qquad H=H_R+\ii H_I,
 \label{eq:quadrature_h}
\end{equation}
where $H_R,H_I$ are real symmetric matrices. They have distinct physical
roles:
\begin{equation}
 \hat h=-\hat D+\ii\hat K,\qquad
 \hat D=-\frac12\hat r^\trans H_R\hat r,\qquad
 \hat K=\frac12\hat r^\trans H_I\hat r.
 \label{eq:damping_unitary}
\end{equation}
When $H_R\leq0$, $\hat D$ is a sum of nonnegative squares of Hermitian
linear combinations of $x,p$. The factor $\ee^{-t\hat D}$ suppresses
large fluctuations, whereas $\ee^{\ii t\hat K}$ is unitary. Their Trotter
product constructs the mixed evolution. Its norm satisfies
\begin{equation}
 \frac{d}{dt}\langle\psi(t)|\psi(t)\rangle
 =-2\langle\psi(t)|\hat D|\psi(t)\rangle\leq0.
 \label{eq:norm_damping}
\end{equation}
Consequently these slices preserve normalizability of every input state.
If the scalar omitted from Eq.~\eqref{eq:quadrature_h} is real, it merely
multiplies this bound by a finite positive exponential.

\subsection{Encoding the same damping in a finite matrix}

The action on the canonical operators follows from one commutator,
\begin{equation}
 [\hat h,\hat r]=-\ii\Omega H\hat r,
 \qquad
 \ee^{-t\hat h}\hat r\,\ee^{t\hat h}=M(t)\hat r,
 \qquad M(t)=\ee^{\ii t\Omega H}.
 \label{eq:propagation}
\end{equation}
The matrix $M$ preserves the canonical commutators,
$M^\trans\Omega M=\Omega$. To measure damping, differentiate
\begin{align}
 \frac{d}{dt}\left[M(t)^\dagger(\ii\Omega)M(t)\right]
 &=M(t)^\dagger(H+H^\dagger)M(t)\nonumber\\
 &=2M(t)^\dagger H_R M(t)\leq0.
 \label{eq:derivative_contraction}
\end{align}
Starting from $M(0)=I$ gives
\begin{equation}
 \boxed{M^\dagger(\ii\Omega)M\leq\ii\Omega.}
 \label{eq:contraction}
\end{equation}
This is the bosonic contraction inequality. The matrix $\ii\Omega$ has
both positive and negative eigenvalues. The word ``contraction'' refers to
this particular comparison of quadratic forms. It does not assert that all
singular values of $M$ are below one.

For example, for $\hat h=-\beta(\hat n+1/2)$ the Nambu column
$\hat\chi=(\hat a,\hat a^+)^\trans$ has propagation
\begin{equation}
 M_\chi=\diag(\ee^{-\beta},\ee^{\beta}),\qquad
 \Sigma=\diag(1,-1),\qquad
 M_\chi^\dagger\Sigma M_\chi-\Sigma
 =\diag(\ee^{-2\beta}-1,1-\ee^{2\beta})<0.
 \label{eq:oscillator_contraction}
\end{equation}
The growing entry in the creation-operator block is required by the
commutator. It is consistent with decay of physical occupation weights.
For $N$ modes the two equivalent forms use
\begin{equation}
 \Sigma=\diag(I_N,-I_N),\qquad
 T=\frac1{\sqrt2}\begin{pmatrix}I&I\\-\ii I&\ii I\end{pmatrix},
 \qquad T\Sigma T^\dagger=\ii\Omega,
 \qquad \hat r=T\hat\chi.
 \label{eq:metric_dictionary}
\end{equation}

\subsection{Why one damping factor can be enough}

A semigroup is simply a set closed under multiplication. Equation
\eqref{eq:contraction} has this property. For $E=\ii\Omega$,
\begin{align}
 E-(M_1M_2)^\dagger E(M_1M_2)
 &=M_2^\dagger(E-M_1^\dagger EM_1)M_2
   +(E-M_2^\dagger EM_2).
 \label{eq:product_contraction}
\end{align}
If both terms are nonnegative, their sum is nonnegative. If one factor has
strict inequality and all the others have nonstrict inequality, the final
product has strict inequality. The propagation matrices are invertible, so
conjugating a positive definite term cannot destroy its strictness.

The factors with $M^\dagger EM=E$ form the \emph{maximal subgroup}: their
inverses obey the same equality. Physically, they are Gaussian unitary
evolutions, apart from the scalar factors already separated, and are
generated here by $H_R=0$. Thus one fully damping factor can be
combined with arbitrarily many such unitary factors. More generally, several
partly damping factors can together damp every direction. Requiring each
slice to have $H_R<0$ is an unnecessarily strong construction rule.

Strict contraction also rules out an eigenvalue on the unit circle. If
$Mv=\lambda v$ with $|\lambda|=1$, then
\begin{equation}
 v^\dagger(E-M^\dagger EM)v=(1-|\lambda|^2)v^\dagger Ev=0,
\end{equation}
contradicting strict positivity. In particular $\lambda=1$, the location of
a bosonic determinant pole, is excluded.

\subsection{Convergence is a property of the final product}
\label{sec:endpoint}

For a nonzero finite-mode Gaussian operator with a finite invertible canonical
propagation matrix, the relation to the actual Fock trace is stronger:
\begin{equation}
 \boxed{\hat G\ \text{is trace class}
 \quad\Longleftrightarrow\quad
 M_\chi^\dagger\Sigma M_\chi<\Sigma.}
 \label{eq:trace_contraction_equivalence}
\end{equation}
The $\Sigma$ in this equation is fixed by the physical relation
$\hat a^+=(\hat a)^\dagger$. Appendix~\ref{app:convergence} proves the
equivalence by Gaussian integrals. It is also part of the usual theory of
Gaussian oscillator operators~\cite{hilgert1989,derezinski2017}.

Equation~\eqref{eq:trace_contraction_equivalence} is a statement about the
\emph{final} operator. Its original generators need not all satisfy
$H_R\leq0$. A simple example is
\begin{equation}
 \ee^{a\hat n}\ee^{-(a+b)\hat n}=\ee^{-b\hat n},
 \qquad a,b>0.
 \label{eq:growth_cancellation}
\end{equation}
Acting from right to left, the first damping creates enough suppression to
permit the later amplification. The final trace is positive and convergent.
The damping inequalities on the separate generators are therefore a
simple sufficient construction of a convergent product. The strict
inequality on the final matrix is the test for that product itself.
We next use the construction to determine the sign.

\section{Fixing the sign by a continuous change of the generators}
\label{sec:trace}
Damping controls the sum over high occupations. The sign can then be
determined by continuously turning the evolution into a simple example
whose trace is known to be positive. During this change, time reversal
keeps the trace real, and contraction keeps it convergent. A finite
Gaussian trace is also nonzero, because it is a nonzero Gaussian
normalization times a convergent Gaussian integral. Appendix
\ref{app:gaussian_trace} derives this last fact explicitly.
The real trace therefore keeps the same sign along the whole path.

\subsection{Starting from a known positive partition function}

The simple example is a set of independent harmonic oscillators.
Choose physical oscillator modes for which
$\hat N_0=\sum_j\hat a_j^+\hat a_j$ is invariant under the common time
reversal. For the actions $\Theta_+\hat a_j\Theta_+^{-1}=\hat a_j$ and
$\Theta_-\hat b_j\Theta_-^{-1}=\hat c_j,\
\Theta_-\hat c_j\Theta_-^{-1}=-\hat b_j$, the ordinary
total occupation already has this property. Appendix
\ref{sec:reference} explains the choice for a general physical time
reversal. This change of physical modes preserves
$\operatorname{Re}H_\ell\leq0$: a nonpositive quadratic form remains
nonpositive under the corresponding real change of $x,p$.
The oscillator partition function is
\begin{equation}
 \Tr\ee^{-\alpha(\hat N_0+N/2)}
 =\prod_{j=1}^N\sum_{n_j=0}^{\infty}
   \ee^{-\alpha(n_j+1/2)}
 =[2\sinh(\alpha/2)]^{-N}>0,\qquad\alpha>0.
 \label{eq:oscillator_partition}
\end{equation}

Now take slices with $\operatorname{Re}H_\ell\leq0$, real scalar terms,
and this common time reversal. Assume the final product is strictly
contracting. To connect its trace to Eq.~\eqref{eq:oscillator_partition},
multiply the original generator by $u$ and add ordinary oscillator
damping with strength $(1-u)\delta$:
\begin{equation}
 \hat h_\ell(u)=u\hat h_\ell-(1-u)\delta
 \sum_j(\hat a_j^+\hat a_j+1/2),
 \qquad 0\leq u\leq1,\quad\delta>0.
 \label{eq:thermal_deformation}
\end{equation}
At $u=0$ all original couplings are absent; at $u=1$ the original
generators are recovered. In these oscillator quadratures,
\begin{equation}
 \operatorname{Re}H_\ell(u)
 =u\operatorname{Re}H_\ell-(1-u)\delta I<0
 \qquad (0\leq u<1).
 \label{eq:deformation_damping}
\end{equation}
Thus every intermediate product has a convergent trace. At $u=1$ the
assumed strict contraction of the complete product supplies convergence
by Eq.~\eqref{eq:trace_contraction_equivalence}.

Write $W(u)=\Tr\prod_\ell\ee^{\hat h_\ell(u)}$. The actual exponentials
vary continuously. Strict contraction supplies a uniform Gaussian decay
bound in a neighborhood of every $u$, including the endpoint $u=1$,
so their Gaussian integrals make $W(u)$ continuous.
Time reversal makes $W(u)$ real. Equation
\eqref{eq:gaussian_integral} ensures that it is never zero. Its initial
value is the oscillator partition function in
Eq.~\eqref{eq:oscillator_partition}, with $\alpha=L\delta$:
\begin{equation}
 W(0)=\left[2\sinh(L\delta/2)\right]^{-N}>0.
 \label{eq:thermal_reference_trace}
\end{equation}
A continuous real number can change sign only by passing through zero.
We have therefore proved
\begin{equation}
 \boxed{\Tr\prod_\ell\ee^{\hat h_\ell}>0.}
 \label{eq:main_positive}
\end{equation}
This proof uses damping, a common time reversal, and an ordinary
Gaussian integral. It applies equally to the two squares of time reversal.

\begin{theorem}[One time reversal in the contraction semigroup]
\label{thm:contraction_trs}
Let
$\hat h_\ell=\frac12\hat r^\trans H_\ell\hat r+\xi_\ell$
be quadratic generators on finitely many bosonic modes, with
$H_\ell^\trans=H_\ell$, $\xi_\ell\in\mathbb R$, and
$\operatorname{Re}H_\ell\leq0$. If all slices preserve a common physical
linear canonical time reversal of square $+I$ or $-I$ on the operators, and
the complete propagation matrix obeys
$M^\dagger(\ii\Omega)M<\ii\Omega$, then the full Fock-space trace exists
and is strictly positive.
\end{theorem}

\subsection{The convergent TRS table and its scope}

Braces conventionally list the squares of common, mutually anticommuting
time reversals. Theorem~\ref{thm:contraction_trs} gives
Table~\ref{tab:trs}. Once one symmetry suffices, adding more symmetries
preserves positivity, regardless of their mutual anticommutation.
\begin{table}[ht]
 \centering
 \begin{tabular}{ll}
 \toprule
 Common time reversals & Convergent contraction construction\\
 \midrule
 None & Positivity is not guaranteed\\
 $\{+\}$ & Positive\\
 $\{-\}$ & Positive\\
 $\{+,+\},\ \{+,-\},\ \{-,-\}$ & All positive\\
 Three or more & All positive\\
 \bottomrule
 \end{tabular}
 \caption{The minimal sufficient classes are $\{+\}$ and $\{-\}$.
 The pairing result follows from Theorem~\ref{thm:contraction_trs}.
 For number conservation, Sec.~\ref{sec:number} requires only convergence
 of the final trace. The extension to amplifying slices is given in
 Appendix~\ref{app:vacuum}.}
 \label{tab:trs}
\end{table}

The table also holds for successive Hermitian and anti-Hermitian slices
acting on a time-reversal-invariant Gaussian vacuum that remains
normalizable through the slices, provided the final trace converges.
Such a factorization can contain
amplification, as in Eq.~\eqref{eq:growth_cancellation}. Appendix
\ref{app:vacuum} supplies this extension and its proof by following a
Gaussian state. The contraction proof above gives the direct route for
damping combined with unitary evolution, which is the construction used
in the examples below.

The role of the path deserves one explicit observation. Multiplying a
Gaussian operator by $-1$ leaves its canonical propagation matrix
unchanged. The matrix inequality alone therefore controls convergence,
while the continuous change of actual exponentials starting from
Eq.~\eqref{eq:thermal_reference_trace} fixes the positive sign.
This is why keeping the physical evolution matters when a trace is
written as a square root of a determinant.

\section{The two standard forms}
\label{sec:standard}

We can now state the connection to \emph{Two standard forms of bosonic
contraction semigroups}, the companion note
\texttt{BosonicContractionSemigroups.md}.
Write time reversal on the Hermitian coordinates as
$\Theta\hat r\Theta^{-1}=R\hat r$. Since it conjugates $\ii$, preservation
of the commutators requires $R^\trans\Omega R=-\Omega$. The conditions on
each quadratic slice are
\begin{equation}
 \boxed{H^\trans=H,\qquad R^\trans\overline H R=H,
 \qquad \operatorname{Re}H\leq0.}
 \label{eq:standard_common}
\end{equation}
They respectively give a quadratic bosonic generator, its physical time
reversal, and damping. In the oscillator coordinates chosen in
Sec.~\ref{sec:trace}, the two
choices of the square of time reversal give the following forms.
A unitary rotation of those modes is enough to obtain them. For square $+I$,
the one-particle space has an orthonormal basis fixed by time reversal:
the combinations of a vector with its time reverse span this real basis.
For square $-I$, a vector and its time reverse are orthogonal, since
$\langle v|\Theta v\rangle=-\langle v|\Theta v\rangle$. Choose one such pair,
then repeat in its orthogonal complement. These are the familiar Kramers
pairs. Both changes of modes preserve physical creation-annihilation
adjoints.

\subsection{Time reversal within each mode: square \texorpdfstring{$+I$}{+I}}

Choose $\Theta\hat x=\hat x$ and $\Theta\hat p=-\hat p$, so
\begin{equation}
 R_+=\begin{pmatrix}I&0\\0&-I\end{pmatrix},\qquad
 H=\begin{pmatrix}H_{xx}&H_{xp}\\H_{xp}^\trans&H_{pp}\end{pmatrix}.
\end{equation}
Substitution into Eq.~\eqref{eq:standard_common} gives
\begin{equation}
 \boxed{H_{xx},H_{pp}\ \text{real},\qquad
 H_{xp}\ \text{purely imaginary},\qquad
 H_{xx}\leq0,\quad H_{pp}\leq0.}
 \label{eq:plus_standard}
\end{equation}
The real $xx$ and $pp$ terms provide damping. The imaginary mixed terms
generate unitary motion that preserves the same time reversal. This is
Theorem 1 of the companion note and the contraction realization of $\{+\}$.

For an explicit example, take one damping factor and one squeezing unitary,
\begin{equation}
 \hat G=\ee^{-\beta(\hat n+1/2)}
 \ee^{\frac{s}{2}(\hat a^2-\hat a^{+2})},
 \qquad \beta>0,\quad s\in\mathbb R.
 \label{eq:plus_example}
\end{equation}
Both generators are real in the occupation basis and preserve $\Theta_+$.
Their trace is
\begin{equation}
 W=\frac{1}{\sqrt{2\cosh\beta\cosh s-2}}>0.
 \label{eq:plus_example_trace}
\end{equation}
At $s=0$ this reduces to the oscillator partition function
$[2\sinh(\beta/2)]^{-1}$. Pairing changes the occupation content.
Along the path in Eq.~\eqref{eq:thermal_deformation}, the trace remains real, convergent,
and nonzero, which fixes the positive square root shown here.

\subsection{Time reversal exchanging two groups: square \texorpdfstring{$-I$}{-I}}

For $N=2m$ modes, use the $b,c$ action in Eq.~\eqref{eq:two_actions} and
order the coordinates as $\hat r=(\hat r_b,\hat r_c)^\trans$, with
$\hat r_b=(\hat x_b,\hat p_b)^\trans$ and likewise for $c$. Define the
momentum-reversal matrix $S=\diag(I_m,-I_m)$. Then
\begin{equation}
 \Theta\hat r_b\Theta^{-1}=S\hat r_c,\qquad
 \Theta\hat r_c\Theta^{-1}=-S\hat r_b.
\end{equation}
Writing $H$ in $b,c$ blocks, Eq.~\eqref{eq:standard_common} becomes
\begin{equation}
 \boxed{H_{cc}=S\overline{H_{bb}}S,\qquad
 H_{bc}=-S H_{bc}^\dagger S,\qquad
 \operatorname{Re}H\leq0.}
 \label{eq:minus_standard}
\end{equation}
The damping inequality applies to the full matrix, including the coupling
between the two groups. The groups need not evolve independently. This is
Theorem 2 of the companion note and the contraction realization of $\{-\}$.

For example, damping of both modes can be combined with the two unitary
generators
\begin{align}
 \hat h_0&=-\beta(\hat n_b+\hat n_c+1),\quad\beta>0,\nonumber\\
 \hat h_1&=\ii\theta(\hat b^+\hat c+\hat c^+\hat b),
 \quad\theta\in\mathbb R,\nonumber\\
 \hat h_2&=\frac12\left(
 \zeta\hat b^{+2}-\overline\zeta\hat b^2
 +\overline\zeta\hat c^{+2}-\zeta\hat c^2\right).
 \label{eq:minus_example}
\end{align}
Applying $\Theta_-$ to each term verifies invariance, including conjugation
of $\ii$ and $\zeta$. The last two generators are anti-Hermitian, and
$\hat h_0$ damps both modes. Every ordered product containing this damping
factor and any sequence of the two unitary factors has a positive convergent
trace. No separation into two independent flavors is used.

\subsection{What the correspondence establishes}

The two forms are the same symmetry distinction that appears in
Theorem~\ref{thm:contraction_trs}. In the original eight-cell enumeration they are
grouped as follows; the cell labels serve as an index into the companion
note:
\begin{center}
\begin{tabular}{lll}
\toprule
Physical TRS & Standard form & Companion-note cells\\
\midrule
$\{+\}$ & Eq.~\eqref{eq:plus_standard} & B02, B05, B17, B22\\
$\{-\}$ & Eq.~\eqref{eq:minus_standard} & B10, B13, B25, B30\\
\bottomrule
\end{tabular}
\end{center}
That enumeration introduces a second anti-linear transformation to encode
the contraction metric through $E=-\overline{R_2}\Omega$.
For the standard metric $E=\ii\Omega$, this gives $R_2=\ii I$.
It is not an additional physical time reversal required of the full
dissipative generator: multiplying Hermitian $x,p$ by $\ii$ would change
their adjoint relation. The full generator needs the one physical symmetry
in Eq.~\eqref{eq:standard_common}.

The two standard forms therefore construct the minimal classes
$\{+\}$ and $\{-\}$ of Table~\ref{tab:trs}. Their generator conditions
keep every multiplication inside the contraction semigroup, while the
connection to independent oscillators fixes the positive sign of the trace.

The companion note additionally identifies eight algebraic descriptions by
complex canonical coordinate changes. When such a change is used, the actual
adjoint relation must be transformed together with the coordinates. The
physical criteria above use Hermitian $x,p$ throughout, so the damping and
Fock-space convergence have their usual oscillator meaning.

\section{Reflection positivity in the same contraction framework}
\label{sec:rp}

Reflection positivity (RP) starts from two groups of bosons. Its physical
idea is to pair each amplitude in one group with its complex conjugate
in the other group. We first recall this construction, then exhibit the
oscillator basis and contraction metric explicitly.

\subsection{The reflected pairs and their common time reversal}

Let $\hat b_j,\hat c_j$, $j=1,\ldots,m$, generate two commuting Fock
spaces. Choose the antilinear reflection
$\theta(\hat b_j)=\hat c_j$,
$\theta(\hat b_j^+)=\hat c_j^+$, and $\theta(\ii)=-\ii$.
Fixed reflection phases can be absorbed into the definitions of the
$c$ modes; this is a unitary mode redefinition preserving their adjoints.
For the present derivation put
\begin{equation}
 \hat z=\binom{\hat b}{\hat b^+},\qquad
 \hat\chi_R=\binom{\hat z}{\theta(\hat z)}
 =(\hat b,\hat b^+,\hat c,\hat c^+)^\trans.
 \label{eq:rp_basis}
\end{equation}
A quadratic RP slice has the form
\begin{equation}
 \hat h=\hat\chi_R^\trans
 \begin{pmatrix}D&F\\F^\trans&\overline D\end{pmatrix}
 \hat\chi_R+\xi,\qquad
 D^\trans=D,\quad F=F^\dagger\geq0,\quad \xi\in\mathbb R.
 \label{eq:rp_generator}
\end{equation}
Each block has dimension $2m$. A nonsymmetric $D$ can be symmetrized:
its antisymmetric part contributes only a scalar, and the two reflected
scalars sum to a real number. The coupling term is
$2\hat z^\trans F\theta(\hat z)$.

Extend the reflection to an antiunitary on the full Fock space by
exchanging the groups in both directions:
\begin{equation}
 \Theta_R\hat b_j\Theta_R^{-1}=\hat c_j,\qquad
 \Theta_R\hat c_j\Theta_R^{-1}=\hat b_j.
 \label{eq:rp_time_reversal}
\end{equation}
It has square $+I$. The two diagonal terms in
Eq.~\eqref{eq:rp_generator} are exchanged. For the mixed term, the
commutation of the two groups and $F_{\nu\mu}=\overline F_{\mu\nu}$
give
\begin{align}
 \Theta_R\!\left[\sum_{\mu\nu}F_{\mu\nu}
       \hat z_\mu\theta(\hat z_\nu)\right]\Theta_R^{-1}
 &=\sum_{\mu\nu}\overline F_{\mu\nu}
       \theta(\hat z_\mu)\hat z_\nu\nonumber\\
 &=\sum_{\mu\nu}F_{\mu\nu}
       \hat z_\mu\theta(\hat z_\nu).
 \label{eq:rp_reality}
\end{align}
Every RP slice therefore has this common physical time reversal.
Hermiticity of $F$ gives the symmetry; its nonnegative eigenvalues
supply the additional reflected-pair positivity.

\subsection{The explicit basis and metric}

In the order of Eq.~\eqref{eq:rp_basis}, the physical adjoint relation
gives
\begin{equation}
 [(\hat\chi_R)_\mu,(\hat\chi_R)_\nu^\dagger]
 =(\eta_R)_{\mu\nu},\qquad
 \boxed{\eta_R=\diag(I_m,-I_m,I_m,-I_m).}
 \label{eq:rp_metric}
\end{equation}
Define $M_R$ by $\hat\chi_R\hat G=\hat G M_R\hat\chi_R$.
For a convergent Gaussian product,
Eq.~\eqref{eq:trace_contraction_equivalence} reads
\begin{equation}
 \boxed{M_R^\dagger\eta_R M_R<\eta_R.}
 \label{eq:rp_contraction}
\end{equation}
The metric is fixed by the physical boson operators.

To put the symmetry itself in the first standard form, introduce ordinary
annihilation operators
\begin{equation}
 \boxed{\hat d_1=\frac{\hat b+\hat c}{\sqrt2},\qquad
        \hat d_2=\frac{\ii(\hat b-\hat c)}{\sqrt2}.}
 \label{eq:rp_d_modes}
\end{equation}
Both obey $\Theta_R\hat d_\alpha\Theta_R^{-1}=\hat d_\alpha$.
For $\hat d_2$, complex conjugation changes the sign of $\ii$, while
the exchange of $b,c$ changes the sign of their difference. The two
signs cancel. These are unitary combinations of annihilation operators,
so the commutators and physical adjoints are preserved.

In full Nambu notation, the transformation is
\begin{equation}
 \hat\chi_d=(\hat d_1,\hat d_2,\hat d_1^+,\hat d_2^+)^\trans
 =S_R\hat\chi_R,\qquad
 S_R=\frac1{\sqrt2}
 \begin{pmatrix}
 I&0&I&0\\
 \ii I&0&-\ii I&0\\
 0&I&0&I\\
 0&-\ii I&0&\ii I
 \end{pmatrix},
 \label{eq:rp_basis_matrix}
\end{equation}
where each $I$ has dimension $m$. Thus
\begin{equation}
 S_R\eta_R S_R^\dagger=\Sigma_d,\qquad
 \Sigma_d=\diag(I_{2m},-I_{2m}),\qquad
 M_d=S_RM_RS_R^\dagger,
 \label{eq:rp_metric_transform}
\end{equation}
and time reversal becomes ordinary coefficient conjugation:
\begin{equation}
 \overline M_d=M_d,\qquad M_d^\trans\Sigma_d M_d<\Sigma_d.
 \label{eq:rp_real_contraction}
\end{equation}
Finally use the Hermitian quadratures
$\hat r_d=(\hat x_{d_1},\hat x_{d_2},
\hat p_{d_1},\hat p_{d_2})^\trans$.
Their metric is $E=\ii\Omega_{2m}$, with
$\Omega_{2m}=\left(\begin{smallmatrix}0&I_{2m}\\-I_{2m}&0\end{smallmatrix}\right)$.
Time reversal fixes their positions and reverses their momenta.
This is precisely the first standard form in
Eq.~\eqref{eq:plus_standard} whenever the slices satisfy the damping
condition. Even if the original factorization includes amplification,
its convergent final matrix still satisfies
Eq.~\eqref{eq:rp_real_contraction}.

\subsection{How RP fixes the positive sign}

Within the damping construction, Theorem~\ref{thm:contraction_trs}
already proves positivity using the oscillator partition function and the symmetry
in Eq.~\eqref{eq:rp_time_reversal}. RP also supplies a direct sign
argument for its own generator cone.

An elementary reflected pair is $\hat R\,\theta(\hat R)$, where
$\hat R$ acts only on the $b$ space. Products of such pairs remain
of the same form because operators on the two spaces commute:
\begin{equation}
 [\hat R\,\theta(\hat R)][\hat Q\,\theta(\hat Q)]
 =(\hat R\hat Q)\theta(\hat R\hat Q).
 \label{eq:rp_product}
\end{equation}
Sums with nonnegative coefficients are therefore closed under
multiplication as well. Decompose
$2F=\sum_\alpha v_\alpha v_\alpha^\dagger$ and set
$\hat A_\alpha=\sum_\mu(v_\alpha)_\mu\hat z_\mu$.
The mixed part of the generator is
$\sum_\alpha\hat A_\alpha\theta(\hat A_\alpha)$.
Its exponential expands in reflected pairs with nonnegative
coefficients. The diagonal terms give
$\ee^{\hat q+\theta(\hat q)}
=\ee^{\hat q}\theta(\ee^{\hat q})$, where
$\hat q=\hat z^\trans D\hat z$.
A Trotter product combines these two pieces. For bounded operators this
constructs limits of finite reflected sums, schematically
\begin{equation}
 \hat G^{(n)}=\ee^{\sum_\ell\xi_\ell}
 \sum_\alpha\hat R_{\alpha,n}\theta(\hat R_{\alpha,n}),
 \label{eq:rp_sum}
\end{equation}
where each sum is finite and positive coefficients are absorbed into
the $\hat R_{\alpha,n}$.
The RP property for an unbounded bosonic construction means that these
finite sums $\hat G^{(n)}$ approach $\hat G$ in every finite-occupation
matrix element.
The series and Trotter argument establish the property at a finite
occupation cutoff. Passing to the original evolution requires the
regularized products to approach its matrix elements.
Convergence of the final trace is then used to remove the cutoff in the
trace, as follows.

The trace of a reflected pair is a modulus square. This can be used
without assuming that the separate factors have finite traces.
Let $P_K$ project onto a finite occupation cutoff in the $b$ space
and use the reflected projector in the $c$ space. With
$\Pi_K=P_K\otimes\theta(P_K)$, the reflected sums give
\begin{equation}
 \Tr(\Pi_K\hat G\Pi_K)
 =\lim_{n\to\infty}\ee^{\sum_\ell\xi_\ell}
 \sum_\alpha
 \left|\Tr_b(P_K\hat R_{\alpha,n}P_K)\right|^2\geq0.
 \label{eq:rp_cutoff_trace}
\end{equation}
The fixed-$K$ trace is a finite sum of matrix elements, so the stated
limit preserves its nonnegative sign.
Since $\hat G$ is trace class, taking $K\to\infty$ yields
$W\geq0$. Its Gaussian trace is nonzero by
Eq.~\eqref{eq:gaussian_integral}, so $W>0$.
In the contraction description, this reflected-pair construction fixes
the positive scalar sign of the operator above the matrix $M_R$.

The two nonnegative conditions have different jobs.
$F\geq0$ pairs reflected amplitudes in Eq.~\eqref{eq:rp_cutoff_trace};
$\operatorname{Re}H\leq0$ suppresses large oscillator fluctuations.
An unrestricted $\hat q+\theta(\hat q)$ can amplify occupations, so
the RP generator condition by itself does not imply slice-by-slice
damping. For every convergent RP product, however, the final
contraction and its positive sign have now both been identified.

\subsection{A paired auxiliary-field example}

The relation becomes particularly concrete for the stable two-flavor
quadratic Hamiltonian
\begin{equation}
 \hat K_0=\sum_{ij}k_{ij}
   (\hat b_i^+\hat b_j+\hat c_i^+\hat c_j)
 -\Delta\sum_i(\hat b_i^+\hat c_i^++\hat b_i\hat c_i),
 \quad k^\trans=k\ \text{real},\quad\Delta\geq0,
 \label{eq:rp_example_hamiltonian}
\end{equation}
with relative-density auxiliary-field factors
\begin{equation}
 \hat T=\ee^{-\Delta\tau\hat K_0},\qquad
 \hat V_\ell=\exp\!\left[
 \ii\sum_i\varphi_{i\ell}(\hat n_{b,i}-\hat n_{c,i})\right],
 \quad \Delta\tau>0,\quad\varphi_{i\ell}\in\mathbb R.
 \label{eq:rp_example_factors}
\end{equation}
The pair term in $-\Delta\tau\hat K_0$ has
$F=(\Delta\tau\Delta/2)I_{2m}\geq0$.
Both factors preserve $\Theta_R$.
If $\lambda_{\min}(k)>\Delta$, all oscillator directions of
$\hat K_0$ have positive energy and $\hat T$ is strictly damping.
Every $\hat V_\ell$ is unitary, so it lies in the maximal subgroup.
Equation~\eqref{eq:product_contraction} proves convergence for any
ordered product containing a $\hat T$ factor and these field factors.

The mode change in Eq.~\eqref{eq:rp_d_modes} gives
\begin{align}
 \hat b_i\hat c_i&=\frac12(\hat d_{1i}^2+\hat d_{2i}^2),\nonumber\\
 \ii\varphi_i(\hat n_{b,i}-\hat n_{c,i})
 &=\varphi_i(\hat d_{1i}^+\hat d_{2i}
             -\hat d_{2i}^+\hat d_{1i}).
 \label{eq:rp_example_transformed}
\end{align}
The pairing is real in each new mode. The field factor is a real
unitary rotation of the two modes. More explicitly, each mode group
has
\begin{equation}
 \hat K_{0,\alpha}
 =\frac12\hat x_{d_\alpha}^\trans(k-\Delta I)\hat x_{d_\alpha}
  +\frac12\hat p_{d_\alpha}^\trans(k+\Delta I)\hat p_{d_\alpha}
  -\frac12\tr k.
 \label{eq:rp_example_quadrature}
\end{equation}
Its negative imaginary-time generator has precisely the nonpositive
$xx,pp$ blocks of Eq.~\eqref{eq:plus_standard}. The field generator
has only imaginary $xp$ coefficients and zero Hermitian part.
Here $\tr k$ is the finite matrix trace.
The path in Eq.~\eqref{eq:thermal_deformation} proves the positive
weight for every field configuration directly.

For nearest-neighbor hopping $t>0$ on the triangular lattice,
$k=tA_\triangle-\mu I$, where $A_\triangle$ is the adjacency matrix,
has its spectrum bounded below by $-3t-\mu$.
The sufficient stability bound becomes
$\mu<-3t-|\Delta|$ after restoring a possible pairing phase.
This stable RP example therefore realizes the first bosonic contraction
semigroup explicitly, with strict damping and unitary field
factors.

\subsection{What is unified by the contraction description}

The metric $\Sigma$ in Nambu coordinates, or $\ii\Omega$ in physical
quadratures, expresses convergence for all the finite Gaussian products
considered here. TRS selects real traces. A continuous connection
to the positive oscillator partition function fixes their sign.
RP supplies another
explicit positive-sign construction through reflected amplitudes.
The two standard forms distinguish square $+I$ from square $-I$ for
the physical time reversal.

For damping and unitary slices, the minimal convergent TRS classes are
therefore $\{+\}$ and $\{-\}$, and the reflected-pair construction
above lies in the first class after Eq.~\eqref{eq:rp_d_modes}.
The symmetry sets $\{-,-\}$, $\{+,-\}$, and $\{+,+,+\}$ are
contained in these convergent sufficient classes. Extending a
determinant formula analytically beyond convergence of the Fock trace
removes the bound that keeps this path free of divergences.
The positivity proof therefore applies to the convergent domain.

The companion note gives local generator conditions that keep every
factor inside either standard contraction semigroup. Convergence of
the final product is more general than these local conditions:
an amplifying factor can be compensated by another factor.
Its final contraction uses the same physical metric; a valid sign
construction, such as RP or the extension in Appendix~\ref{app:vacuum},
specifies the positive operator trace.

\appendix
\section{Why a finite Gaussian trace cannot pass through zero}
\label{app:gaussian_trace}

Use coherent states
\begin{equation}
 |z\rangle=\ee^{-|z|^2/2}
 \exp\!\left(\sum_jz_j\hat a_j^+\right)|0\rangle,
 \qquad
 \int\frac{d^{2N}z}{\pi^N}|z\rangle\langle z|=I.
 \label{eq:coherent_states}
\end{equation}
The matrix element of a quadratic evolution between coherent states is a
Gaussian. For some symmetric matrices $X,V$ and a matrix $Y$, its unnormalized
form is
\begin{equation}
 \langle0|\ee^{\bar z^\trans\hat a}\hat G
 \ee^{w^\trans\hat a^+}|0\rangle
 =c\exp\!\left[
 \frac12\bar z^\trans X\bar z+\bar z^\trans Yw
 +\frac12w^\trans Vw\right].
 \label{eq:gaussian_kernel}
\end{equation}
Here $c=\langle0|\hat G|0\rangle\ne0$ for a nonzero finite Gaussian
evolution. The Gaussian form follows by applying the linear propagation
relations to the coherent-state kernel; Sec.~\ref{app:kernel} derives
its coefficients. In particular, the kernel is a nonzero constant times
an exponential, so its value at $z=w=0$ cannot vanish.

Set $w=z$ and include the coherent-state normalization. Writing the real and
imaginary parts of $z$ as a real vector $u$ turns the trace into an ordinary
Gaussian integral. Its absolute convergence follows from the convergence
already assumed for $\hat G$. To see this explicitly, write a singular-value
decomposition $\hat G=\sum_j s_j|u_j\rangle\langle v_j|$, with normalized
states and $\sum_j s_j<\infty$. Coherent-state completeness and the
Cauchy--Schwarz inequality give
\begin{equation}
 \int\frac{d^{2N}z}{\pi^N}|\langle z|\hat G|z\rangle|
 \leq\sum_j s_j
 \left[\int\frac{d^{2N}z}{\pi^N}|\langle z|u_j\rangle|^2
       \int\frac{d^{2N}z}{\pi^N}|\langle v_j|z\rangle|^2\right]^{1/2}
 =\sum_j s_j.
 \label{eq:coherent_absolute}
\end{equation}
Thus the real part of the quadratic decay matrix in this integral must be
positive definite. The result is
\begin{equation}
 \Tr\hat G
 =c\int\frac{d^{2N}u}{(2\pi)^N}
       \exp\!\left(-\frac12u^\trans\mathcal A u\right)
 =\frac{c}{\sqrt{\det\mathcal A}},
 \qquad \operatorname{Re}\mathcal A>0.
 \label{eq:gaussian_integral}
\end{equation}
The normalization uses $z=(u_x+\ii u_p)/\sqrt2$. The positive real part is
the convergence condition for this Gaussian integral. It also makes
$\mathcal A$ nonsingular. One can evaluate the integral by continuously
changing $\mathcal A$ from $I$ to its final value while keeping its real
part positive. Hence its value is finite and nonzero. No sign choice from a
standalone determinant is needed: the convergent integral defines the value.

\subsection{The Gaussian kernel and its propagation matrix}
\label{app:kernel}

This calculation supplies the finite-matrix connection used in
Secs.~\ref{sec:trace} and \ref{sec:endpoint}. Only linear propagation
relations and coherent-state differentiation are needed.

In the Nambu basis define the propagation by
\begin{equation}
 \hat\chi\hat G=\hat G M_\chi\hat\chi,
 \qquad
 M_\chi=\begin{pmatrix}A&B\\C&D\end{pmatrix}.
 \label{eq:kernel_intertwine}
\end{equation}
The letters $A,B,C,D$ here denote propagation blocks, independently of
the RP generator blocks in Sec.~\ref{sec:rp}.
With $z,w$ independent complex variables, define
\begin{equation}
 \mathcal K(z,\bar w)=
 \langle0|\ee^{z^\trans\hat a}\hat G
 \ee^{\bar w^\trans\hat a^+}|0\rangle.
 \label{eq:appendix_kernel}
\end{equation}
This is an analytic, nonzero kernel when $\hat G$ is a nonzero bounded
operator.
Creation acts by multiplication and annihilation by differentiation on
these unnormalized coherent states. The lower block of
Eq.~\eqref{eq:kernel_intertwine} therefore gives
\begin{equation}
 z\mathcal K=C\bar w\mathcal K+D\partial_{\bar w}\mathcal K.
\end{equation}
If $v^\trans D=0$, this equation would require
$v^\trans(z-C\bar w)\mathcal K=0$ for all independent $z,w$.
The kernel is analytic and nonzero, so $v=0$ and $D$ is invertible.
Substituting its inverse and then using the upper propagation block gives
\begin{align}
 \partial_{\bar w}\mathcal K
   &=(D^{-1}z-D^{-1}C\bar w)\mathcal K,\nonumber\\
 \partial_z\mathcal K
   &=(BD^{-1}z+D^{-\trans}\bar w)\mathcal K.
\end{align}
Here the canonical relations give
$A-BD^{-1}C=D^{-\trans}$. These first-order equations fix the kernel
up to a constant $c$, with the Gaussian form
\begin{equation}
 \mathcal S=
 \begin{pmatrix}
 BD^{-1}&D^{-\trans}\\
 D^{-1}&-D^{-1}C
 \end{pmatrix},
 \qquad
 \mathcal K=c\exp\!\left[\frac12
 (z,\bar w)^\trans\mathcal S(z,\bar w)\right].
 \label{eq:kernel_S}
\end{equation}
The coefficients in Eq.~\eqref{eq:gaussian_kernel} are therefore
$X=BD^{-1}$, $Y=D^{-\trans}$, and $V=-D^{-1}C$.
The constant $c$ cannot be zero, since that would make the whole kernel
and the operator zero. Its value is supplied by the actual evolution,
for example by Eq.~\eqref{eq:c_evolution}. The propagation matrix
determines the shape of the Gaussian kernel.

\section{Deriving the equivalence between convergence and strict contraction}
\label{app:convergence}

For a Gaussian kernel the convergence test is another Gaussian integral.
The sum of the squared singular values of the operator is
\begin{equation}
 \Tr(\hat G^\dagger\hat G)
 =|c|^2\int\frac{d^{4N}v}{\pi^{2N}}
 \exp\!\left[-v^\dagger v+\operatorname{Re}(v^\trans\mathcal Sv)\right],
 \label{eq:kernel_norm_integral}
\end{equation}
where $v=(z,\bar w)$ ranges over $2N$ independent complex variables.
A unitary change of these variables diagonalizes the symmetric quadratic
form as $\sum_js_jv_j^2$, with real $s_j\geq0$. These $s_j$ are the
singular values of $\mathcal S$. In the real and imaginary parts of each
$v_j$, the Gaussian decay coefficients are $1-s_j$ and $1+s_j$.
Therefore
\begin{equation}
 \Tr(\hat G^\dagger\hat G)<\infty
 \quad\Longleftrightarrow\quad
 I-\mathcal S^\dagger\mathcal S>0.
 \label{eq:kernel_norm_test}
\end{equation}
Trace class implies the left-hand condition, since a summable sequence of
nonnegative singular values is also square summable.

Conversely, in this finite-mode Gaussian setting square summability already
implies trace class. If $\|\mathcal S\|<1$, choose a small $\epsilon>0$
with $\ee^{2\epsilon}\|\mathcal S\|<1$. Multiplying every occupation
matrix element on the left and right by $\ee^{\epsilon\hat N_0}$ scales
the kernel matrix to $\ee^{2\epsilon}\mathcal S$, which still defines a
bounded operator $\hat G_\epsilon$. Then
\begin{equation}
 \hat G=\ee^{-\epsilon\hat N_0}\hat G_\epsilon
        \ee^{-\epsilon\hat N_0}.
 \label{eq:traceclass_sandwich}
\end{equation}
Both outside factors have square-summable singular values. The singular-value
Cauchy--Schwarz inequality bounds the sum of singular values of this product
by $\|\ee^{-\epsilon\hat N_0}\|_2^2\|\hat G_\epsilon\|$, which is
finite. Here $\|T\|_2^2=\Tr(T^\dagger T)$. This proves trace class.

It remains to express Eq.~\eqref{eq:kernel_norm_test} using $M_\chi$.
Define the invertible block matrix
\begin{equation}
 L_0=\begin{pmatrix}0&I\\D^{-1}&-D^{-1}C\end{pmatrix}.
\end{equation}
The canonical relations give $A-BD^{-1}C=D^{-\trans}$. Consequently
\begin{equation}
 L_0=\begin{pmatrix}0&I\\Y^\trans&V\end{pmatrix},
 \qquad
 M_\chi L_0=\begin{pmatrix}X&Y\\I&0\end{pmatrix}.
\end{equation}
These two block matrices make the remaining calculation explicit:
\begin{align}
 L_0^\dagger\Sigma L_0
 &=\begin{pmatrix}
 -\overline Y Y^\trans&-\overline Y V\\
 -V^\dagger Y^\trans&I-V^\dagger V
 \end{pmatrix},\nonumber\\
 (M_\chi L_0)^\dagger\Sigma(M_\chi L_0)
 &=\begin{pmatrix}
 X^\dagger X-I&X^\dagger Y\\
 Y^\dagger X&Y^\dagger Y
 \end{pmatrix}.
\end{align}
Subtracting them, and using
$\mathcal S=\left(\begin{smallmatrix}X&Y\\Y^\trans&V\end{smallmatrix}\right)$,
gives
\begin{equation}
 L_0^\dagger(\Sigma-M_\chi^\dagger\Sigma M_\chi)L_0
 =I-\mathcal S^\dagger\mathcal S.
 \label{eq:congruence_identity}
\end{equation}
Since $L_0$ is invertible, the two sides are positive definite together.
Equations~\eqref{eq:kernel_norm_test}--\eqref{eq:congruence_identity}
prove Eq.~\eqref{eq:trace_contraction_equivalence}.

For completeness, strict contraction itself guarantees invertibility of
$D$ before the kernel is constructed. Testing
$\Sigma-M_\chi^\dagger\Sigma M_\chi>0$ on a vector $(0,w)^\trans$
gives $\|Dw\|^2>\|w\|^2+\|Bw\|^2$ for nonzero $w$. Thus $D$ has no
null vector. Equation~\eqref{eq:kernel_S} then constructs the convergent
Gaussian kernel for the given matrix, with its overall scalar fixed by the
physical evolution. Finite products of Gaussian exponentials are the
nondegenerate operators under discussion; an exact vacuum projector is a
limiting object rather than one with a finite invertible $M_\chi$.

\section{An extension allowing amplifying slices}
\label{app:vacuum}
\subsection{Choosing oscillator coordinates that respect time reversal}
\label{sec:reference}

To use the independent-oscillator partition function in
Sec.~\ref{sec:trace}, we need a total occupation that is invariant under
time reversal. For the familiar actions
\begin{equation}
 \Theta_+\hat a_j\Theta_+^{-1}=\hat a_j,\qquad
 \Theta_-\hat b_j\Theta_-^{-1}=\hat c_j,\quad
 \Theta_-\hat c_j\Theta_-^{-1}=-\hat b_j,
 \label{eq:two_actions}
\end{equation}
the total occupation is already invariant. The second action requires an
even number of modes. Its square is $-I$ on the canonical operators and
particle parity on Fock space.

A general physical canonical time reversal can be put in this form after
choosing suitable oscillator modes. To obtain a positive reference
Hamiltonian, start from any positive confining oscillator $\hat H_0$ and
take
\begin{equation}
 \hat H_{\rm ref}=\hat H_0+\Theta\hat H_0\Theta^{-1}.
 \label{eq:reference_hamiltonian}
\end{equation}
Both terms are positive confining quadratics, and the sum is invariant
because applying $\Theta$ twice leaves a quadratic unchanged. Diagonalize
this stable oscillator Hamiltonian into normal modes. Time reversal
mixes annihilation operators only with annihilation operators at the same
positive frequency, as also follows from their commutator with
$\hat H_{\rm ref}$. It therefore preserves the total normal-mode
occupation $\hat N_0=\sum_j\hat a_j^+\hat a_j$.
The mode frequencies can be replaced by one common positive frequency
while preserving the symmetry. In these physical modes,
$\hat N_0+N/2=\hat r^\trans\hat r/2$ is the desired reference.
The change of modes preserves physical adjoints and the sign of the
quadratic damping condition.

\subsection{A normalizable reference state}
\label{sec:vacuum}

The continuous path in Sec.~\ref{sec:trace} gives a direct proof
for damping and unitary slices. This appendix treats a broader evolution
that may include amplification. The additional object is one normalizable
Gaussian reference state followed through the slices. Write the general
normally ordered quadratic generator as
\begin{equation}
 \hat h=\xi+(\hat a^+)^\trans K\hat a
 +\frac12(\hat a^+)^\trans P\hat a^+
 +\frac12\hat a^\trans Q\hat a,
 \qquad P^\trans=P,\quad Q^\trans=Q.
 \label{eq:normal_generator}
\end{equation}
The columns $\hat a$ and $\hat a^+$ contain $N$ annihilation and creation
operators, with $[\hat a_i,\hat a_j^+]=\delta_{ij}$. The matrices $P,Q,K$ need
not obey Hermiticity relations.

Use the reference modes of Sec.~\ref{sec:reference} and choose the phase
of their vacuum so that $\Theta|0\rangle=|0\rangle$.
The unique ground state of Eq.~\eqref{eq:reference_hamiltonian} has
this property. Time reversal acts on its annihilation operators as
in Eq.~\eqref{eq:trs_a}.

The state remains centered, meaning that it has no displacement
$\langle\hat a_j\rangle$. A centered pure Gaussian state can be written as
\begin{equation}
 |\psi\rangle=c\,|Z\rangle,
 \qquad
 |Z\rangle=\exp\!\left[\frac12(\hat a^+)^\trans Z\hat a^+\right]|0\rangle,
 \qquad Z^\trans=Z.
 \label{eq:squeezed_state}
\end{equation}
The exponential creates pairs. Its constant term is one, so
$c=\langle0|\psi\rangle$ is exactly the vacuum amplitude, including its
sign or phase. This convention removes any freedom to attach a changing
phase to $|Z\rangle$.

For one mode,
\begin{equation}
 \langle Z|Z\rangle
 =\sum_{n=0}^{\infty}\frac{(2n)!}{2^{2n}(n!)^2}|Z|^{2n}
 =\frac{1}{\sqrt{1-|Z|^2}}.
 \label{eq:squeezed_norm_one}
\end{equation}
The state is normalizable exactly when $|Z|<1$. A unitary rotation of modes
reduces a complex symmetric $Z$ to independent pair amplitudes $z_j$; in
several modes,
\begin{equation}
 \langle Z|Z\rangle=\prod_j(1-|z_j|^2)^{-1/2},
 \qquad \|Z\|<1.
 \label{eq:squeezed_norm}
\end{equation}
Here $|z_j|$ are the singular values of $Z$, and $\|Z\|$ is the largest
one. This elementary condition is how normalizability will enter the proof.

\subsection{Following the vacuum amplitude through a time slice}
\label{sec:evolution}

Let a slice act for $0\leq t\leq1$, with
$\partial_t|\psi(t)\rangle=\hat h|\psi(t)\rangle$. Quadratic evolution
preserves the Gaussian form, so insert Eq.~\eqref{eq:squeezed_state} and
determine $c(t),Z(t)$. The commutation relation first gives
\begin{align}
 \hat a_i|Z\rangle&=\sum_j Z_{ij}\hat a_j^+|Z\rangle,
 \label{eq:one_annihilator}\\
 \hat a_i\hat a_j|Z\rangle
 &=Z_{ij}|Z\rangle+
   \sum_{kl}Z_{ik}Z_{jl}\hat a_k^+\hat a_l^+|Z\rangle.
 \label{eq:two_annihilators}
\end{align}
The first term in Eq.~\eqref{eq:two_annihilators} is the contribution that
can change the vacuum amplitude. Comparing the terms with zero and two
creation operators on the two sides of the evolution equation yields
\begin{align}
 \frac{\dot c}{c}&=\xi+\frac12\tr(QZ),
 \label{eq:c_evolution}\\
 \dot Z&=P+KZ+ZK^\trans+ZQZ.
 \label{eq:z_evolution}
\end{align}
The symbol $\tr$ denotes a finite matrix trace; $\Tr$ denotes a many-body
trace. Equation~\eqref{eq:z_evolution} determines the pair content. For the
sign, only Eq.~\eqref{eq:c_evolution} is needed.

Because $\Theta$ preserves $\hat h$ and the initial state, the evolving
state is invariant. Acting with Eq.~\eqref{eq:trs_a} on its pair exponential
and on the annihilation-pair term in $\hat h$ gives
\begin{equation}
 Z=U^\dagger\overline Z\,\overline U,
 \qquad Q=U^\trans\overline Q\,U,
 \qquad \xi=\overline\xi.
 \label{eq:qz_reality}
\end{equation}
These relations make the coefficient in Eq.~\eqref{eq:c_evolution} real.
The nontrivial check is
\begin{equation}
 \overline{\tr(QZ)}
 =\tr\!\left(\overline U Q U^\dagger UZ U^\trans\right)
 =\tr(QZ),
 \label{eq:trace_qz_real}
\end{equation}
where $U^\trans\overline U=I$ was used in the last step. The sign of
$\overline U U$ has not entered.

Starting from the vacuum, $c(0)=1$. As long as the Gaussian state remains
normalizable during the evolution, $Z$ stays finite and
\begin{equation}
 c(t)=\exp\!\left\{\int_0^t
 \left[\xi+\frac12\tr(QZ(\tau))\right]d\tau\right\}>0.
 \label{eq:c_positive}
\end{equation}
There is no sign ambiguity in this equation: it is the exponential of an
ordinary real integral. Apply this argument to the slices in their order of
action, from the rightmost factor in Eq.~\eqref{eq:weight} to the leftmost.
The positive amplitude at the end of one slice is the starting amplitude for
the next. Therefore
\begin{equation}
 \langle0|\hat G|0\rangle>0.
 \label{eq:vacuum_positive}
\end{equation}

\subsection{When does the required state evolution exist?}
\label{sec:existence}

The state condition in this proof is concrete: the Gaussian obtained at
every intermediate time must have finite norm, equivalently $\|Z(t)\|<1$.
There are several directly useful ways to ensure it.

A Hermitian slice has this property whenever its complete action on the
input state exists. To see this, resolve that state into spectral components
of $\hat h$, of eigenvalue $\lambda$. For $0\leq t\leq1$,
\begin{equation}
 \|\ee^{t\hat h}\psi\|^2
 =\int\ee^{2t\lambda}\,d\mu_\psi(\lambda)
 \leq\int(1+\ee^{2\lambda})\,d\mu_\psi(\lambda)
 =\|\psi\|^2+\|\ee^{\hat h}\psi\|^2.
 \label{eq:hermitian_slice}
\end{equation}
The measure $d\mu_\psi$ is the squared amplitude of the state in this
spectral resolution. The inequality proves that all intermediate states
exist, even for an amplifying Hermitian slice such as $+a\hat n$.
An anti-Hermitian slice generates a unitary evolution and preserves the norm.
Thus a product of Hermitian and anti-Hermitian slices needs no separate
intermediate-state assumption when the stated successive actions exist.

Another useful choice is a slice with a nonpositive Hermitian quadratic part.
Its evolution decreases the norm, up to a finite scalar factor. This allows
mixed, non-Hermitian slices and is the contraction construction in
Sec.~\ref{sec:contraction}. More generally, one can use
Eq.~\eqref{eq:z_evolution} and check $\|Z(t)\|<1$ for a chosen reference
vacuum. This version places no sign restriction on the instantaneous
Hermitian part.

This check follows the state during each time slice. If a singular value of
$Z(t)$ reaches one, Eq.~\eqref{eq:squeezed_norm} diverges, and that vacuum
trajectory cannot be used in the proof. The constructions above ensure
finite norm throughout the slice.

\subsection{From the vacuum amplitude to the full trace}
\label{sec:filter}

A positive vacuum amplitude is only one term of the occupation sum. We now
show why it fixes the sign of the whole Gaussian trace. The physical operation
is to suppress excited oscillator states more and more strongly. Define
\begin{equation}
 \hat N_0=\sum_j\hat a_j^+\hat a_j,
 \qquad F(s)=\Tr\!\left(\ee^{-s\hat N_0}\hat G\right),
 \qquad s\geq0.
 \label{eq:thermal_filter}
\end{equation}
The reference modes are those of Sec.~\ref{sec:vacuum}, so $\hat N_0$ is
time-reversal invariant. At $s=0$ this is the required weight. At large $s$
only the vacuum survives.

Three properties give the result. First, the added factor has norm at most
one, so it preserves convergence of the trace. In the reference occupation
basis,
\begin{equation}
 F(s)=\sum_{\bb n}\ee^{-s\sum_jn_j}
 \langle\bb n|\hat G|\bb n\rangle,
 \qquad
 \lim_{s\to\infty}F(s)=\langle0|\hat G|0\rangle.
 \label{eq:filter_limit}
\end{equation}
Absolute convergence permits this limit and makes $F(s)$ continuous.
Second, time reversal preserves both factors, so $F(s)$ is real for all $s$.
Third, a convergent trace of a nonzero centered Gaussian evolution cannot
vanish. This was proved by the coherent-state integral in
Appendix~\ref{app:gaussian_trace}.

The same reasoning applies to $\ee^{-s\hat N_0}\hat G$ at every finite
$s\geq0$. Its vacuum coefficient is still $c$. Thus $F(s)$ is a real,
continuous function that never reaches zero and approaches the positive
number in Eq.~\eqref{eq:vacuum_positive}. It follows that
\begin{equation}
 \boxed{W=F(0)>0.}
 \label{eq:filter_positive}
\end{equation}
This identifies the role of convergence in the pairing problem. It allows
the excited states to be suppressed continuously, without a divergent
Gaussian integral interrupting the connection to the vacuum.

\subsection{The extended sufficient condition}

\begin{theorem}[One physical time reversal]
\label{thm:one_trs}
Let $\hat G=\prod_\ell\ee^{\hat h_\ell}$ be an actual ordered quadratic
evolution on finitely many bosonic modes with a convergent full Fock-space
trace. Suppose all generators preserve a common physical antiunitary
$\Theta$, with square $+I$ or $-I$ on the single-particle operators.
If a time-reversal-invariant Gaussian vacuum has normalizable evolution
through these slices, then $\Tr\hat G>0$.

In particular, this state condition holds for successive Hermitian and
anti-Hermitian slices whose actions exist, and for the dissipative slices of
Sec.~\ref{sec:contraction}. It can also be checked directly through
Eq.~\eqref{eq:z_evolution}.
\end{theorem}

\section{Checking the one-mode example}

For Eq.~\eqref{eq:plus_example}, put $q=\ee^{-\beta}$ and $v=\tanh s$.
The propagation and the blocks of the kernel are
\begin{equation}
 M_\chi=
 \begin{pmatrix}q&0\\0&q^{-1}\end{pmatrix}
 \begin{pmatrix}\cosh s&-\sinh s\\-\sinh s&\cosh s\end{pmatrix},
 \quad X=-q^2v,\quad Y=\frac{q}{\cosh s},\quad V=v.
\end{equation}
The vacuum amplitude is $c=\sqrt{q/\cosh s}$. Substituting the diagonal
kernel into the two-real-variable integral in
Eq.~\eqref{eq:gaussian_integral} gives
\begin{equation}
 \frac{W}{c}=\frac{1}{\sqrt{(1-Y)^2-XV}},
 \qquad
 W^{-2}=q^{-1}\cosh s+q\cosh s-2.
\end{equation}
This is Eq.~\eqref{eq:plus_example_trace}. The direct integral agrees
with the positive sign fixed by the path in
Eq.~\eqref{eq:thermal_deformation}.

\begin{thebibliography}{9}
\bibitem{hilgert1989}
J.~Hilgert, A note on Howe's oscillator semigroup,
\emph{Ann. Inst. Fourier} \textbf{39}, 663--688 (1989),
\href{https://doi.org/10.5802/aif.1182}{doi:10.5802/aif.1182}.
\bibitem{derezinski2017}
J.~Derezi\'nski and M.~Karczmarczyk, Quantization of Gaussians (2017),
\href{https://arxiv.org/abs/1701.07297}{arXiv:1701.07297}.
\end{thebibliography}
\end{document}
